\section{Conclusion}
\label{sec:conclusion}

We asked whether safety tuning in open-weight chat models creates surface-level blind spots in the way that gradient masking does in adversarial robustness. Across 14 models from 0.5B to 32.8B parameters and four distribution shifts applied to the same harmful requests, the answer depends on the shift. Models are not keyword matchers: refusal correlates 1.5--4$\times$ more with semantic harm than with trigger words at every size, and removing trigger words costs a moderate 16.3\pp of refusal (11.8\pp when rated harm is held fixed) with only a 3.5\pp rise in harmful completions. Multilingual non-refusal in small models is mostly incoherence, not harmful help. The exploitable blind spots are framing blind spots: one paragraph of fiction or a no-rules persona raises compliance to 40--93\% in seven of the eleven safety-tuned models, and these gaps are family-specific and not closed by scale.

{\bf The key takeaway is that safety tuning generalises across words but not across frames.} Future work should validate the judge with human annotation, especially for non-English replies; use longer generations to measure how much actionable detail fictional frames leak; test system prompts and sampling; probe activations to test whether framing suppresses a harm representation or the refusal direction \cite{arditi2024refusal}; and extend the study to API models and to optimised paraphrases such as PAIR \cite{chao2023pair}.
